\documentclass[11pt,a4paper]{article}
\usepackage[T1]{fontenc}
\usepackage[utf8]{inputenc}
\usepackage{amsmath,amssymb,amsthm}
\usepackage[margin=1in]{geometry}
\usepackage[colorlinks=true,linkcolor=blue,urlcolor=blue]{hyperref}
\theoremstyle{plain}
\newtheorem{theorem}{Theorem}
\newtheorem{lemma}[theorem]{Lemma}
\newtheorem{proposition}[theorem]{Proposition}
\newtheorem{corollary}[theorem]{Corollary}
\theoremstyle{definition}
\newtheorem{remark}[theorem]{Remark}

\title{The empirical recurrence for OEIS A264109, proved by transfer matrix}
\author{Adrian Perez Fontelles\\ \small Independent researcher}
\date{1 September 2026}
\begin{document}
\maketitle

\begin{abstract}
OEIS A264109 counts the arrangements of $0,1,\dots,LW-1$ in a grid of $W=7$ rows in which
every value moves from its home cell by an index change on a short list, and it carries an
empirical recurrence of order $25$ contributed by R.~H.~Hardin. It is true, and it is
decidable rather than empirical. Such an arrangement is a perfect matching between cells and
values, and because a value's row moves by at most $R=1$, the matching splits into
independent decisions row by row: after the cells of rows $0,\dots,i$ are filled,
every value whose home row is at most $i-R$ has been used, and only $2R$ home rows are
partly used. Carrying those as bitmasks makes the count a walk on $S=16384$ vertices, so it is
$C$-finite and the conjectured recurrence is decided by a finite exact computation.
\end{abstract}

\noindent\small 2020 Mathematics Subject Classification. 05A05, 05A15, 15A15.\normalsize

\section{The conjecture}

OEIS A264109 is ``Number of (n+1)X(6+1) arrays of permutations of 0..n*7+6 with each element having index change +-(.,.) 0,0 0,2 or 1,1.''. It has offset $1$ and begins
\[
1936, 167464, 15437041, 1399892079, 127528266321, 11613782698504, 1057423575802500, 96278613956842527,\ \dots
\]
The entry states, as a conjecture contributed by its author and never marked settled:
\begin{quote}
Empirical: a(n) = 87*a(n-1) +32068*a(n-3) +180766*a(n-4) -3715198*a(n-5) -16639004*a(n-6) +149707392*a(n-7) +542256817*a(n-8) +224301365*a(n-9) +4280490864*a(n-10) -2613943512*a(n-11) +12541107972*a(n-12) -12056948868*a(n-13) +9621158616*a(n-14) -5946708960*a(n-15) -2779571151*a(n-16) +6286349025*a(n-17) -2468635920*a(n-18) -364811148*a(n-19) +148489470*a(n-20) -3874014*a(n-21) +2657124*a(n-22) +16767*a(n-24) +243*a(n-25)
\end{quote}
As of the ``Last modified'' line on the live entry (Nov 10 2015 11:26:05, revision 6) this is still
recorded as empirical, and nothing on the entry records it as proved.

\section{The object is a matching}

The grid has $W=7$ rows across and holds each of $0,1,\dots,LW-1$ exactly once. The value
$v$ has a HOME cell, namely its place in the plain row-major filling, and putting $v$ in cell
$(i,j)$ changes its index by $(i-h_i,\,j-h_j)$ where $(h_i,h_j)$ is that home cell. The entry
allows the index changes
\[
(i-h_i,\,j-h_j)\;\in\;\{(-1,-1),\ (0,-2),\ (0,0),\ (0,2),\ (1,1)\},
\]
the list being the entry's own, read with the sign convention its wording uses. An admissible
arrangement is therefore exactly a perfect matching in the bipartite graph joining each cell to
the values it may hold.

\section{The matching is a walk}

Let $R=1$ be the largest $|i-h_i|$ on the list. A value whose home row is $h$ can only be
placed in rows $h-R,\dots,h+R$. So once the cells of rows $0,\dots,i$ are filled, every
value with home row at most $i-R$ has been used up, and the only home rows that can be
partly used are $i-R+1,\dots,i+R$ --- exactly $2R$ of them. Take as the state the tuple of
$2R$ bitmasks recording which values of those home rows are already placed. One step fills
the cells of one further row, drawing each of its $W$ values from a home row within $R$
and at an allowed column offset, and the home row that closes at that step must come out
full. Nonexistent home rows --- before the first and after the last --- are carried as
full, which makes the starting and accepting states the same one.

Hence $a$ is a walk count on $S=16384$ vertices,
\[
a(n)\;=\;\iota^{\!\top}M^{\,n+1}\tau ,
\]
the exponent fixed by the entry's own shape and checked against its published terms; in
particular $a$ is $C$-finite.

\section{The proof}

\begin{lemma}\label{lem:crit}
Let $q(t)=t^{25}-\sum_i c_it^{25-i}$ be the characteristic polynomial of the
conjectured recurrence. The residual $a(n)-\sum_ic_ia(n-i)$ equals
$\iota^{\!\top}M^{\,j}q(M)\tau$ for the corresponding walk index $j$, so the recurrence
holds from some point on if and only if $\iota^{\!\top}M^{j}q(M)\tau=0$ for all large $j$;
and it is enough to examine $j<S$.
\end{lemma}

\begin{proof}
Factor $M^{j}$ out of $M^{j+25}-\sum_ic_iM^{j+25-i}$. For the bound, $M^{S}$
is an integer combination of $I,M,\dots,M^{S-1}$ by Cayley--Hamilton, so the numbers
$u_j=\iota^{\!\top}M^{j}q(M)\tau$ obey the monic recurrence given by the characteristic
polynomial of $M$; $S$ consecutive zeros force every later one, and the last nonzero $u_j$
pins the threshold exactly.
\end{proof}

\begin{theorem}
\[
a(n)\;=\;87\,a(n-1) + 32068\,a(n-3) + 180766\,a(n-4) - 3715198\,a(n-5) - 16639004\,a(n-6) + 149707392\,a(n-7) + 542256817\,a(n-8) + 224301365\,a(n-9) + 4280490864\,a(n-10) - 2613943512\,a(n-11) + 12541107972\,a(n-12) - 12056948868\,a(n-13) + 9621158616\,a(n-14) - 5946708960\,a(n-15) - 2779571151\,a(n-16) + 6286349025\,a(n-17) - 2468635920\,a(n-18) - 364811148\,a(n-19) + 148489470\,a(n-20) - 3874014\,a(n-21) + 2657124\,a(n-22) + 16767\,a(n-24) + 243\,a(n-25)
\]
for every $n>25$.
\end{theorem}

\begin{proof}
The vector $w=q(M)\tau$ was formed by $25$ matrix--vector products in exact integer
arithmetic and $\iota^{\!\top}M^{j}w$ evaluated for $j=0,1,\dots$, again exactly; those
integers vanish from the index corresponding to $n=25+1$ onwards and the run continued
until the working vector was identically zero, settling every larger $j$ at once.
\end{proof}

\section{Verification}

Three checks, each independent of the algebra above.

First, before any of this was built the reading of the entry's wording was tested against the
entries themselves by brute force: enumerating all $720$ arrangements of a $2\times3$ grid
gives $20$, the first term of A263960, and the $3\times2$ grid gives $9$, the first term of its
transposed companion A263966. Second, the transfer model reproduces all $12$ terms this
entry publishes, exactly, in integer arithmetic. That check earned its keep: the entries use
two different sign conventions, ``$(\pm,\pm)$ $a,b$'' meaning each coordinate signed
independently and ``$\pm(.,.)$ $a,b$'' meaning the PAIR signed, where the second coordinate may
itself be written negative. Reading the second as the first silently drops a displacement, and
thirteen of the first twenty-three entries tested disagreed with their own data until the two
were separated.

Third, the annihilation test of Lemma~\ref{lem:crit} was carried out in exact integer
arithmetic on the whole vector rather than on a sampled prefix.

\begin{thebibliography}{9}
\bibitem{oeis} The OEIS Foundation, \emph{The On-Line Encyclopedia of Integer
Sequences}, \url{https://oeis.org}, 2026. Sequence A264109.
\bibitem{stanley} R.~P.~Stanley, \emph{Enumerative Combinatorics}, Volume 1, 2nd ed.,
Cambridge University Press, 2011. (Transfer-matrix method, Section 4.7.)
\bibitem{lp} L.~Lov\'asz and M.~D.~Plummer, \emph{Matching Theory}, AMS Chelsea, 2009.
\bibitem{horn} R.~A.~Horn and C.~R.~Johnson, \emph{Matrix Analysis}, 2nd ed., Cambridge
University Press, 2013.
\end{thebibliography}

\end{document}
